\documentclass[a4paper,14pt]{extarticle}
\usepackage[utf8]{inputenc}
\usepackage[T2A]{fontenc}
\usepackage[russian]{babel}
\usepackage{geometry}
\usepackage{graphicx}
\usepackage{array}
\usepackage{longtable}
\usepackage{titlesec}
\usepackage{hyperref}
\usepackage{setspace}
\usepackage{indentfirst}
\usepackage{fancyhdr}

\geometry{left=3cm, right=1.5cm, top=2cm, bottom=2cm}
\onehalfspacing
\setlength{\parindent}{1.25cm}

\pagestyle{fancy}
\fancyhf{}
\fancyfoot[C]{\thepage}
\renewcommand{\headrulewidth}{0pt}

\titleformat{\section}[block]{\normalfont\bfseries\centering}{}{0pt}{}
\titleformat{\subsection}[block]{\normalfont\bfseries}{}{0pt}{}

\begin{document}

% ========== ТИТУЛЬНЫЙ ЛИСТ ==========
\begin{titlepage}
\begin{center}

\small
МИНИСТЕРСТВО НАУКИ И ВЫСШЕГО ОБРАЗОВАНИЯ РОССИЙСКОЙ ФЕДЕРАЦИИ

\smallskip
Федеральное государственное автономное образовательное учреждение высшего образования

\smallskip
\textbf{<<Дальневосточный федеральный университет>>}

\textbf{(ДВФУ)}

\vspace{5mm}
\noindent\rule{\textwidth}{2pt}

\vspace{10mm}
{\large\textbf{ИНСТИТУТ МАТЕМАТИКИ И КОМПЬЮТЕРНЫХ ТЕХНОЛОГИЙ}}

{\large\textbf{(ШКОЛА)}}

\vspace{5mm}
\textbf{Департамент математического и компьютерного моделирования}

\vspace{15mm}
{\large\textbf{Лабораторная работа №1}}

\vspace{5mm}
по теме: <<Анализ рисков информационной безопасности>>, вариант 10

\end{center}

\vspace{30mm}

\begin{flushright}
Выполнил студент\\
гр. Б9123-02.03.01сцт\\
Избосарова Жасмина Олимжоновна
\end{flushright}

\vfill

\begin{center}
г. Владивосток\\
2026
\end{center}

\end{titlepage}

% ========== СОДЕРЖАНИЕ ==========
\tableofcontents
\newpage

% ========== ЗАДАНИЕ ==========
\section*{Задание}
\addcontentsline{toc}{section}{Задание}

В данной работе рассматривается дизайнерская фирма (вариант 10). Задание выполняется по методичке.

\begin{center}
\begin{tabular}{|c|p{13cm}|}
\hline
\textbf{№} & \textbf{Формулировка задания} \\
\hline
1 & Загрузите ГОСТ Р ИСО/МЭК ТО 27005-2010 <<МЕТОДЫ И СРЕДСТВА ОБЕСПЕЧЕНИЯ БЕЗОПАСНОСТИ>>. Часть 3 <<Методы менеджмента безопасности информационных технологий>>. \\
\hline
2 & Ознакомьтесь с Приложениями C, D и E ГОСТа. \\
\hline
3 & Выберите три различных информационных актива организации (см. вариант). \\
\hline
4 & Из Приложения D ГОСТа подберите три конкретных уязвимости системы защиты указанных информационных активов. \\
\hline
5 & Пользуясь Приложением С ГОСТа напишите три угрозы, реализация которых возможна пока в системе не устранены названные в пункте 4 уязвимости. \\
\hline
6 & Пользуясь одним из методов предложенных в Приложении Е ГОСТа произведите оценку рисков информационной безопасности. \\
\hline
7 & Оценку ценности информационного актива производить на основании возможных потерь для организации в случае реализации угрозы. \\
\hline
\end{tabular}
\end{center}

% ========== 1. ОБОСНОВАНИЕ ВЫБОРА ==========
\section{Обоснование выбора информационных активов организации}

Рассматривается небольшая дизайнерская фирма, которая занимается разработкой фирменного стиля, полиграфии, веб-дизайна и рекламных материалов для заказчиков. В штате около 15 человек: дизайнеры, менеджеры проектов и бухгалтер. Для работы нужны три вещи: файлы с проектами (макеты, исходники), рабочие компьютеры дизайнеров и электронная почта для общения с клиентами. Именно эти три актива и были выбраны~--- они покрывают и творческую часть, и рабочий процесс, и коммуникации.

\begin{center}
\begin{tabular}{|c|p{4cm}|p{8.5cm}|}
\hline
\textbf{№} & \textbf{Информационный актив} & \textbf{Назначение и обоснование выбора} \\
\hline
1 & Файловое хранилище проектов & Здесь лежат все рабочие файлы: макеты в Figma/Photoshop/Illustrator, исходники логотипов, вёрстка, шрифты. Это главный результат работы фирмы. Если файлы пропадут или утекут к конкурентам~--- фирма потеряет и работу, и доверие клиентов. \\
\hline
2 & Рабочие станции дизайнеров & Компьютеры (обычно Mac или мощные ПК), на которых стоят лицензионные программы: Adobe Creative Cloud, Figma, CorelDRAW. На них ведётся вся работа. Поломка или заражение вирусом означает простой. \\
\hline
3 & Корпоративная электронная почта & Через почту идёт переписка с заказчиками: техзадания, правки, утверждение макетов, отправка готовых файлов. Также через неё приходят договоры и счета. Взлом почты~--- это и утечка информации, и потеря клиентов. \\
\hline
\end{tabular}
\end{center}

Все три актива связаны друг с другом: дизайнер работает на рабочей станции, сохраняет файлы в хранилище, а результат отправляет клиенту по почте. Если что-то из этого выйдет из строя или будет скомпрометировано, пострадает вся цепочка.

% ========== 2. ОЦЕНКА ЦЕННОСТИ ==========
\section{Оценка ценности информационных активов}

Ценность каждого актива оценивается по тому, какие потери понесёт фирма, если с ним что-то случится. Фирма небольшая, поэтому суммы не такие большие, как в банке, но для маленькой компании они ощутимы. Используется шкала от 1 до 5 по примеру 2 из Приложения Е.2.2: 1~--- минимальные потери, 5~--- очень серьёзные.

\begin{center}
\begin{tabular}{|c|p{3.2cm}|p{4.5cm}|c|p{3.5cm}|}
\hline
\textbf{№} & \textbf{Информационный актив} & \textbf{Возможные потери при реализации угрозы} & \textbf{V (1--5)} & \textbf{Обоснование} \\
\hline
1 & Файловое хранилище проектов & Около 1,5 млн руб.: переделка утраченных проектов, срыв сроков, штрафы по договорам, потеря клиентов из-за утечки их материалов. & 5 & Это основной продукт работы фирмы. Потеря файлов~--- потеря всего. \\
\hline
2 & Рабочие станции дизайнеров & Около 500 тыс. руб.: простой сотрудников, покупка/переустановка ПО, восстановление локальных файлов. & 3 & Неприятно, но восстановить можно за несколько дней. Лицензии привязаны к аккаунтам, а не к железу. \\
\hline
3 & Корпоративная электронная почта & Около 800 тыс. руб.: утечка переписки с клиентами, потеря контактов, рассылка спама от имени фирмы, репутационный ущерб. & 4 & Почта~--- лицо компании. Взлом сильно бьёт по репутации. \\
\hline
\end{tabular}
\end{center}

% ========== 3. УЯЗВИМОСТИ ==========
\section{Уязвимости системы защиты информации}

Уязвимости взяты из Приложения D, таблица D.1. Для каждого актива подобрана та уязвимость, которая реально может встретиться в небольшой дизайнерской фирме, где обычно нет выделенного сисадмина и вопросы безопасности решаются <<по ходу дела>>.

\begin{center}
\begin{tabular}{|c|p{3cm}|p{2.2cm}|p{4cm}|p{3cm}|}
\hline
\textbf{№} & \textbf{Информационный актив} & \textbf{Вид} & \textbf{Конкретная уязвимость из Приложения D} & \textbf{Связанная примером угроза} \\
\hline
1 & Файловое хранилище проектов & Программные средства & Отсутствие резервных копий & Искажение данных \\
\hline
2 & Рабочие станции дизайнеров & Персонал & Ненадлежащее использование программных и аппаратных средств & Ошибка в использовании \\
\hline
3 & Корпоративная электронная почта & Сеть & Передача паролей в открытом виде & Перехват информации \\
\hline
\end{tabular}
\end{center}

Уязвимости выбраны не случайно. В маленькой фирме часто нет бэкапов (<<да ладно, ничего не случится>>), сотрудники ставят что попало на рабочие компьютеры, а пароли от почты бывают простыми или передаются в мессенджерах открытым текстом.

% ========== 4. УГРОЗЫ ИБ ==========
\section{Угрозы ИБ}

Угрозы подобраны из Приложения С, таблица С.1. <<У>>~--- умышленная угроза, <<С>>~--- случайная. Для каждой пары <<уязвимость--актив>> описан конкретный сценарий: что может случиться и как.

\begin{center}
\begin{tabular}{|c|p{2.8cm}|p{2.5cm}|c|p{5cm}|}
\hline
\textbf{№} & \textbf{Уязвимость} & \textbf{Угроза ИБ из Приложения С} & \textbf{Тип} & \textbf{Сценарий возможного инцидента} \\
\hline
1 & Отсутствие резервных копий & Искажение данных & С & Дизайнер случайно удаляет папку с проектом крупного заказчика. Бэкапов нет, корзина очищена. Проект приходится делать заново, срок сорван, клиент уходит. \\
\hline
2 & Ненадлежащее использование ПО и аппаратных средств & Ошибка в использовании & С & Сотрудник скачивает <<бесплатный>> плагин для Photoshop с сомнительного сайта. Вместе с плагином устанавливается вредоносная программа, которая шифрует файлы на компьютере. Работа останавливается, пока комп не переустановят. \\
\hline
3 & Передача паролей в открытом виде & Перехват информации & У & Менеджер отправляет пароль от корпоративной почты новому сотруднику в обычном сообщении в Telegram. Злоумышленник, получивший доступ к одному из аккаунтов, видит пароль и заходит в почту фирмы. Читает переписку, скачивает файлы клиентов. \\
\hline
\end{tabular}
\end{center}

Эти три сценария вполне реалистичны для небольшой фирмы. Дальше на их основе считается риск.

% ========== 5. ОЦЕНКА РИСКОВ ==========
\section{Оценка рисков}

Используется подход из Приложения Е, пример Е.2.2: риск считается как произведение ценности актива на вероятность реализации угрозы. Формула: $R = V \times P$.

Шкала вероятности:

\begin{center}
\begin{tabular}{|c|p{3.5cm}|p{7cm}|}
\hline
\textbf{P} & \textbf{Степень вероятности} & \textbf{Что это значит} \\
\hline
1 & Очень низкая & Маловероятно, нужно редкое стечение обстоятельств. \\
\hline
2 & Низкая & Может случиться, но нечасто. \\
\hline
3 & Средняя & Вполне реальный сценарий. \\
\hline
4 & Высокая & Скорее всего случится при текущих условиях. \\
\hline
5 & Очень высокая & Случается регулярно или почти неизбежно. \\
\hline
\end{tabular}
\end{center}

\vspace{3mm}

Расчёт рисков:

\begin{center}
\begin{tabular}{|c|p{2.8cm}|p{2.8cm}|p{2.5cm}|c|c|c|c|}
\hline
\textbf{№} & \textbf{Актив} & \textbf{Уязвимость} & \textbf{Угроза} & \textbf{V} & \textbf{P} & \textbf{R} & \textbf{Ранг} \\
\hline
1 & Файловое хранилище проектов & Отсутствие резервных копий & Искажение данных & 5 & 4 & 20 & 1 \\
\hline
2 & Корпоративная электронная почта & Передача паролей в открытом виде & Перехват информации & 4 & 3 & 12 & 2 \\
\hline
3 & Рабочие станции дизайнеров & Ненадлежащее использование ПО & Ошибка в использовании & 3 & 3 & 9 & 3 \\
\hline
\end{tabular}
\end{center}

\vspace{3mm}

Итоговое ранжирование:

\begin{center}
\begin{tabular}{|c|p{9cm}|c|p{3cm}|}
\hline
\textbf{Приоритет} & \textbf{Сценарий} & \textbf{R} & \textbf{Итоговая оценка} \\
\hline
1 & Потеря файлов проектов из-за отсутствия бэкапов & 20 & Очень высокий риск \\
\hline
2 & Перехват пароля и взлом корпоративной почты & 12 & Высокий риск \\
\hline
3 & Заражение рабочей станции из-за скачивания сомнительного ПО & 9 & Средний риск \\
\hline
\end{tabular}
\end{center}

Самый опасный сценарий~--- потеря проектных файлов ($R = 5 \times 4 = 20$). Вероятность поставлена 4, потому что в маленьких фирмах бэкапы часто не делаются вообще. На втором месте~--- перехват пароля от почты ($R = 4 \times 3 = 12$). На третьем~--- заражение рабочей станции ($R = 3 \times 3 = 9$).

% ========== 6. ВЫВОДЫ ==========
\section{Выводы}

В работе были рассмотрены три информационных актива дизайнерской фирмы: файловое хранилище проектов, рабочие станции дизайнеров и корпоративная электронная почта.

Для каждого актива определена ценность от 3 до 5 по пятибалльной шкале. Самым ценным оказалось хранилище проектов~--- потеря файлов для дизайнерской фирмы фактически означает потерю результатов работы.

Из Приложения D подобраны три уязвимости: отсутствие бэкапов, небрежное обращение с ПО и передача паролей в открытом виде. Каждой уязвимости сопоставлена угроза из Приложения С.

По методу Е.2.2 получены меры риска: 20, 12 и 9. Главный риск~--- потеря данных без возможности восстановления. Чтобы его снизить, нужно в первую очередь настроить автоматическое резервное копирование (хотя бы в облако). Кроме того, стоит ввести правило не передавать пароли через мессенджеры и не скачивать программы из непроверенных источников.

% ========== 7. ЛИТЕРАТУРА ==========
\section{Литература}

\begin{enumerate}
\item ГОСТ Р ИСО/МЭК 27005-2010, Приложения C, D и E.
\end{enumerate}

\end{document}
